\section{Experimental Method}
\label{sec:setup}

\subsection{One variable}
\label{sec:determinism}

A migrated run and its control must differ only in the state that crossed. Three mechanisms enforce this. The batch
used at step $t$ is drawn from a generator seeded by a hash of the run seed, the stream and $t$, so the data order is a
function of the step index and no loader state exists to be restored incorrectly. Dropout masks are seeded the same
way. The learning rate is a function of $t$. Training, the destination's warm-up passes and held-out evaluation use
three separate streams, so the warm-up passes neither consume nor repeat training batches. Deterministic kernels are
enforced in the framework. Under these conditions a migration with the \arm{full} strategy reproduces the control's
loss at every later step bit for bit, on the same device; the test suite asserts this, and \arm{full} appears in the
tables as a row of exact zeros. This is the bitwise standard EasyScale sets for elastic
training~\cite{li2023easyscale}, used here as an instrument: any gap in any other row was caused by the missing
state.

Each control saves a full snapshot at every migration point. A migration cell loads the snapshot into a source
process, packs the strategy's payload, writes it, reads it into a freshly built destination, rebuilds what is missing
and trains for $W = \horizon$ further steps. Nothing but the bundle passes from one to the other.

\subsection{Metrics}

Loss is evaluated every \evalEvery{} steps on a fixed held-out set of \evalSamples{} sequences. Training loss is too
noisy at these sizes and is confounded with which batch lands on which step. Let $L_{\mathrm{mig}}$ and
$L_{\mathrm{ctrl}}$ be the held-out losses of a migrated run and of its own control, with the same seed, and
$\Delta(j) = L_{\mathrm{mig}}(T+j) - L_{\mathrm{ctrl}}(T+j)$.

\textbf{Seed band.} For each model size and migration point, $\eps$ is the 95th percentile of
$|L_a(t) - L_b(t)|$ over all pairs of control seeds $a \ne b$ and all evaluations in $(T, T+W]$. It measures how far
apart two runs of the same job are for no reason at all. A migrated run with $\Delta \le \eps$ is no further from its
control than another seed would be. With \nSeedsMax{} seeds there are three such pairs at each point, so the band is
a coarse estimate; across the sweep it lies between \epsMin{} and \epsMax\mnats.

\textbf{Peak excess loss} is $\max_j \Delta(j)$, reported in thousandths of a nat (mnats) and as a multiple of $\eps$.
It is the main quantity in the tables because it is defined for every run and has a direct reading: a ratio below one
means the run never left the band.

\textbf{Steps to recovery} (StR) is the first offset $k$ such that $\Delta(j) \le \eps$ for all $j \in [k, k+50]$.
One lucky crossing does not count. A run with no such $k$ inside the window is censored and reported as $>W$; it is
never assigned the window length. Tables give the number of runs with $\mathrm{StR} = 0$, which never left the band,
the median StR of those that left and returned, and the number censored.

\textbf{Regret} is $\sum_j \max(0, \Delta(j))$ over the window, in nat-steps. It separates a deep, brief transient from
a shallow, long one.

\textbf{Steps lost} converts the training-loss gap at the end of the window into the number of steps by which the
migrated run trails its control, read off the control's smoothed loss curve, plus any warm-up passes spent. It is the
unit a scheduler wants, and it is noisy: applied to pairs of controls it gives values of \bandLamMin{} to
\bandLamMax{} steps with no migration at all. We report it with that caveat and use it only in the cost model.

\textbf{Uncertainty.} The migration points of one seed share a control and are not independent. Intervals are
therefore 95\% percentile intervals from \bootDraws{} bootstrap resamples of whole seeds.

\subsection{Workloads and hardware}

The main workload is a decoder-only transformer trained on character-level text~\cite{karpathy2022nanogpt} at three
sizes, \paramsXS{}M, \paramsS{}M and \paramsM{}M parameters, for \maxSteps{} steps with AdamW, $\beta_1 = 0.9$,
$\beta_2 = 0.99$, weight decay 0.1, gradient clipping at 1, a \warmupSteps-step warm-up and cosine decay, batch 64 and
context 256, under bfloat16 autocast with single-precision master weights and moments. Each size is run with
\nSeedsMax{} seeds and migrated at steps \forkList{}, that is, 10\%, 30\%, 60\% and 85\% of the schedule. A second sweep repeats nine of the strategies on the middle size at
$\beta_2 = 0.95$ and $0.999$, with migrations at 10\% and 60\% (Section~\ref{sec:beta}), and a twelve-layer model of \largeParams{}M parameters is
run with five strategies at the same two points (Section~\ref{sec:large}).

All but one experiment ran on one laptop: an Intel Core Ultra 9 185H with 32\,GB of memory, an NVIDIA RTX 4050 Laptop GPU with
6\,GB, and the processor's integrated Arc GPU, under Windows 11 with PyTorch 2.14 (CUDA 12.6) and, for the integrated
GPU, PyTorch 2.15. The one exception is the remote resume of Section~\ref{sec:device}, which ran on an NVIDIA T4 in a
hosted notebook under Linux with PyTorch 2.11. Every run writes a manifest with the commit, the resolved configuration, the device, the measured
payload size and the median step time.
